%==============================================================================
%  Section 6 --- Proof Outline (NEW)                                [agent A2]
%
%  Sources: v8 thesis_v8.tex:482-560 (Neumann motivation, thm:neumann, sketch),
%           v8:969-978 (lem:C-bern, promoted here as lem:pernode),
%           v7 thesis_v7.tex:755-849 (proof flowchart, reconciled to v8 values).
%  Renames: thm:neumann -> lem:neumann (demoted: not a result of this paper)
%           lem:C-bern  -> lem:pernode (statement promoted; proof stays in app:comp3b)
%  [POLISH-A]: flowchart removed, essays cut to one paragraph/stage -- see 17-polish-a.md
%
%  [Reviewer Note (v9 refine pass): subsections 5.1-5.5 removed.  A proof sketch
%  is the one place in a paper where momentum matters most -- the reader is
%  being carried through an argument, not consulting a manual -- and five
%  headings across two pages announced each step twice, once in the roadmap
%  sentence and once as a title.  The stages now run as continuous prose;
%  lem:neumann, lem:pernode and the proof environment are untouched.  Internal
%  \Cref{sec:outline-*} self-references became prose ("above", "below"), and
%  appendix-C's two references to sec:outline-synth now point at sec:outline.]
%
%  [Reviewer Note (v9 outline pass): the section promised three stages but
%  stated only two lemmas, and the two paragraphs between the opening and the
%  proof re-derived the margin, the gap and the noise bound in prose -- so the
%  same algebra appeared twice on one page, once informally and once inside a
%  lemma.  Three changes.  (i) The STATEMENT of lem:B-bernstein is promoted here
%  from app:comp2, giving Stage 1 a lemma of its own; three stages now carry
%  three lemmas, and its proof stays in the appendix.  (ii) The motivational
%  prose is rewritten at the conceptual altitude of v8's "The Neumann trick"
%  paragraph (v8:495-517), which explains WHY the device is decisive without
%  competing with the lemma that follows it; every formula it used to carry is
%  now stated once, in a lemma.  The standalone motivational paragraph that
%  briefly sat between the opening and Stage 1 was then removed as well: it
%  argued method choice rather than this proof (nothing in Steps 1-4 depended on
%  it), it landed after the opening's handoff sentence so the promised Stage 1
%  was 250 words late, and it pre-empted Stage 3's own punchline.  Its three
%  load-bearing pieces were redistributed to where they do work -- the
%  l-infinity/l-2 mismatch now opens Stage 2 as the reason that stage exists, the
%  attribution and the lemma-not-theorem note ride on lem:neumann's pointer, and
%  the sqrt(m) saving plus the app:compare/rem:E-compare pointers close Stage 3.
%  The same argument is made at length in appendix-E and in intro_prev.tex.
%  (iii) Steps 1-3 of the proof are one clause
%  each, naming only what each stage hands to the next; Step 4 is the single
%  place in the section where algebra is done.  The Step 1-4 labels are
%  retained: appendix-C cites "Step 2" and "Step 4", and app:dist cites "the
%  same three stages and synthesis".]
%==============================================================================
\section{Proof Outline}\label{sec:outline}
The proof of \Cref{thm:main-sim} compares two quantities evaluated at the same
node: the entry-wise error $|\vvh_i-\vv_i|$ and the entry-wise margin $|\vv_i|$.
It runs in three stages --- population geometry fixes the margin and noise
budget, a resolvent expansion reduces the error at each node to a scalar, and a
per-node Bernstein bound controls that scalar --- after which imposing
\eqref{eq:entrywise} yields the rate \eqref{eq:main_rate}. Each stage is stated
as a lemma; the proofs are in \Cref{app:comp1,app:comp2,app:comp3}.

\emph{Stage 1 (population geometry and noise budget).} The same imbalance $\eta$
sets both sides of \eqref{eq:entrywise} at once --- it fixes the margin and the
spectral gap through the piecewise-constant Fiedler vector of
\Cref{lem:A-decomp} and the gap bound of \Cref{lem:gap}, kept positive by
\Cref{asm:margin}, while on the other side matrix Bernstein applied to the
sampling mask bounds the noise the gap has to absorb.

\begin{lemma}[Laplacian spectral error via matrix Bernstein]\label{lem:B-bernstein}
Under \Cref{asm:clock,asm:cfbm} and the observation model
\eqref{eq:sampling_estimator}, there are constants $\tilde C,C_0$
such that if $p\ge C_0\,\tfrac{\log m}{m}$ then \hp
\[
   \opnorm{\EL}\le\tilde C\sqrt{\frac{m\log m}{p}}.
\]
\end{lemma}
\noindent The proof is in \Cref{app:comp2}. Compared against the gap bound of
\Cref{lem:gap}, this is what certifies the strong-signal regime
$\opnorm{\EL}<\Dlam$ that the next stage assumes. Under imbalance the margin and
the gap are extremal at different nodes, which is what makes the final synthesis
$\eta$-dependent.

\emph{Stage 2 (entry-wise reduction).} Since \eqref{eq:entrywise} is an
$\ell_\infty$ event while Davis--Kahan controls only $\ell_2$, and the sole
dimension-free transfer between them, $\norm{\cdot}_\infty\le\norm{\cdot}_2$,
charges every node as though it were the worst one, the reduction has to be
coordinate-wise: in the regime of Stage 1 a resolvent (Neumann) expansion of
$\vvh$ in $\EL$ trades the operator norm for a single scalar per coordinate, at
the price of a remainder that is second order in the noise-to-gap ratio.

\begin{lemma}[Neumann decomposition of the Fiedler vector]\label{lem:neumann}
Let $L$ be symmetric with eigenpairs $\{(\lambda_\ell,v^{(\ell)})\}$, let
$\hat L=L+\EL$ be a symmetric perturbation, and write $\vvh=\vv+w$. If
$\opnorm{\EL}<\Dlam$, then, with $\Pi$ the spectral projector onto the complement
of $\vv$, $w$ admits the exact resolvent expansion
$w=\sum_{k\ge1}(-1)^k(\Dlam^{-1}\Pi\EL)^k\vv$, whose first-order term is
\begin{equation}\label{eq:neumann_first_order}
   w^{(1)}=\sum_{\ell\ne2}\frac{-\langle v^{(\ell)},\EL\vv\rangle}
        {\lambda_\ell-\hat\lambda_2}\,v^{(\ell)},
\end{equation}
so that entrywise $w_i=-(\EL\vv)_i/\Dlam+R_i$ with the remainder of second order,
$\norm{R}_\infty=O\!\bigl(\opnorm{\EL}^2/\Dlam^2\bigr)$.
\end{lemma}
\noindent The expansion is due to Eldridge et al.\
\cite{eldridge2017unperturbed}, Fan et al.\ \cite{fan2018eigenvector} and Abbe
et al.\ \cite{abbe2020entrywise} rather than to this paper, which is why it is
stated as a lemma; the decomposition and its first-order bound are proved in
\Cref{app:comp3a}, and the remainder $R$ separately in \Cref{app:comp3c}
(\Cref{lem:C-remainder}), its discard being what confines the argument to the
strong-signal regime of Stage 1.

\emph{Stage 3 (per-node concentration).} The surviving scalar $(\EL\vv)_i$ is a
sum of independent, mean-zero increments, each carrying the factor $\vv_i-\vv_j$.
Every within-clan increment therefore vanishes \emph{exactly} --- this is a
cancellation, not a bound, and it rests on the clock (\Cref{asm:clock}) making
$\vv$ exactly piecewise constant, so it is the one step that would degrade if the
cross-clan block were only approximately constant. Only cross-clan pairs survive,
Bernstein's inequality applies node by node, and since this is the only place the
sampling rate $p$ enters the final condition, we state the bound in the main text.

\begin{lemma}[Per-node concentration]\label{lem:pernode}
With scale $R_B=\bze(1+\eta)/(p\sqrt{m\eta})$, for each $i\in B$ and $\delta>0$,
with probability $\ge1-\delta$,
$|(\EL\vv)_i|\le2\sqrt{2\sigma_B^2\ln(2/\delta)}+\tfrac23R_B\ln(2/\delta)$. In the
variance-dominated regime ($p=\Theta(\log m/m)$), a union bound over all $m$ nodes
at $\delta=\delta_0/m$ gives, \hp\ uniformly over $i$,
\begin{equation}\label{eq:concentration-bound}
   |(\EL\vv)_i|\lesssim\sqrt{\frac{8(1+\eta)\bze^2\log(2m/\delta_0)}{\eta\,p}}.
\end{equation}
\end{lemma}
\noindent The increment representation, the cancellation of the within-clan
pairs, the per-node variance $\sigma_B^2$ and the proof are in
\Cref{app:comp3b}. This is the step that realises the $\sqrt m$ saving that the
Davis--Kahan route forfeits: the variance of the increment sum is smaller than
$\opnorm{\EL}^2$ by a factor $m$, and it is this per-coordinate saving that turns
\eqref{eq:entrywise} into the rate \eqref{eq:main_rate}; \Cref{app:compare} runs
the two routes side by side and \Cref{rem:E-compare} reports the ratio. The three
stages now combine.

\begin{proof}[Proof of \Cref{thm:main-sim}]
\emph{Step 1 (population geometry).} \Cref{lem:A-decomp,lem:gap} supply the
right-hand side of \eqref{eq:entrywise} as $\min_i|\vv_i|=c_B$ and a spectral gap
$\Dlam$ that is positive under \Cref{asm:margin}, and \Cref{lem:B-bernstein}
places $\opnorm{\EL}$ below it.

\emph{Step 2 (entry-wise reduction).} \Cref{lem:neumann} is therefore applicable
and gives $w_i=-(\EL\vv)_i/\Dlam+R_i$; discarding $R$ as lower order against
$c_B$ reduces \eqref{eq:entrywise} to the per-node condition
$|(\EL\vv)_i|<\Dlam|\vv_i|$ holding at every node $i$.

\emph{Step 3 (concentration).} \Cref{lem:pernode} bounds the left-hand side of
that condition uniformly over the $m$ nodes.

\emph{Step 4 (synthesis).} The per-node condition is tightest on clan $B$, since
$\eta\ge1$. Squaring $\sqrt{8\sigma_B^2\log m}<\Dlam c_B$ and substituting
$\sigma_B^2$, $c_B^2=1/(m\eta)$ and the gap bound
$\Dlam\ge m_{\min}(\rho-\eta\Sout^{\max})$, the factor $\eta$ cancels and solving
for $p$ gives exactly \eqref{eq:main_rate}
(\Cref{prop:C-bind} and \Cref{cor:C-cond}). Under \eqref{eq:main_rate} every node satisfies
$|w_i|<|\vv_i|$, hence $\vvh_i$ and $\vv_i$ share a sign; since $\vv$ is
sign-constant on each clan (\Cref{lem:A-decomp}), the zero-threshold partition of
$\vvh$ is $(A,B)$ exactly.
\end{proof}

\Cref{thm:main-dist} states the counterpart of \Cref{thm:main-sim} for the centered
distance operator. Its proof follows the same four steps, with the population
geometry of Step~1 and the variance load of Step~3 replaced by their distance
analogues; \Cref{lem:neumann}
is applied unchanged, being a statement about an arbitrary symmetric perturbation
of an arbitrary symmetric matrix. \Cref{app:dist} carries the details.
